\section{Models, Datasets, and Scale}
\label{sec:setup}

\paragraph{Checkpoints.}
We evaluate 21 instruction-tuned checkpoints from nine families (Table~\ref{tab:cf-families}; full list in Table~\ref{tab:cf-models}): Qwen2.5-VL 3B/7B/32B/72B, Qwen3-VL 4B/8B/32B, InternVL3.5 8B/14B/38B, Gemma~3 4B/12B/27B, MedGemma 4B/27B, Lingshu 7B/32B, LLaVA-1.5 7B/13B, LLaVA-Med~1.5 7B, and Llama~3.2~Vision 11B. Three families are medical fine-tunes (MedGemma, Lingshu, LLaVA-Med); the rest are general-purpose. The families span five vision-tower designs and three consumer designs (token insertion after a merger, token insertion after a projector, and cross-attention), and consumed-token counts from 144 to 4{,}096 per image. Every checkpoint runs in bf16 at a pinned revision with its native processor settings and a fixed image resolution per family.

\paragraph{Datasets and concepts.}
NIH ChestX-ray14 \citep{wang2017chestxray} and CheXpert Plus provide chest radiographs with report-derived labels; we evaluate six findings on each (NIH: Effusion, Atelectasis, Pneumothorax, Cardiomegaly, Mass, Nodule; CheXpert: Effusion, Atelectasis, Pneumothorax, Cardiomegaly, Consolidation, Edema). COCO provides the natural-image control with six object categories (person, dog, car, chair, bottle, bicycle) as binary presence labels. Each dataset has one fixed cohort chosen by seeded hashes of patient (or image) identity: 20{,}000 training rows, 400 calibration rows, 600 test rows, and 16 preflight rows, patient-disjoint and shared by every checkpoint, one frontal study per patient on the chest sets.

\paragraph{Labels.} Concept labels are the report-derived labels shipped with each dataset (NIH ChestX-ray14 finding labels; the CheXpert labeler output distributed with CheXpert Plus; COCO instance annotations). A row is positive when the label is 1, negative when it is 0, and \emph{unknown} when the labeler returns uncertain or blank; unknown rows are excluded from probe fitting, selectivity, and answer AUROC, and are written and scored like every other row. Cohorts are fixed manifests drawn with recorded seeds, patient-disjoint between training and test and disjoint from every seed-study cohort; per-concept counts are in Table~\ref{tab:cf-prevalence}.

\paragraph{Questions and templates.}
Each concept is asked as a yes/no question under six templates that cross two wordings (\emph{is there} versus \emph{does the image show}) with three answer interfaces (yes/no, A-present, B-present). The primary grade uses the \emph{is}/yes-no template; the other five feed the wording and mapping contrasts. The main write dose is $\alpha=+0.25$; the dose module adds $\alpha\in\{-0.5,-0.25,-0.1,+0.1,+0.5\}$.

\paragraph{Scale.}
A block is one checkpoint on one dataset. Each chest block scores the $6\times6$ write matrix, the sham, and 119 random directions on 600 test rows (457{,}200 outcomes), plus calibration, dose, refit, connector-locus, and template modules; COCO blocks run the full module set and CheXpert blocks the write matrix and calibration. The campaign comprises 61 blocks, 342 scored concept cells, about 30 million scored answers, and roughly 3{,}000 A100-hours. Nineteen checkpoints are complete on all three datasets; Llama~3.2~Vision 11B is graded on its eligible templates; Qwen2.5-VL-72B is run on NIH only, and no checkpoint above 72B is run (Appendix~\ref{app:campaign}).

\paragraph{Registration.} The protocol was frozen before any block of the grid was scored: the loci, the projection and probe settings, the dose, the reference family, the six templates, the module set, the seven gates, and the three decision rules are those of the registered protocol file and its code commit, and the seed study's crossover and confirmation analyses were registered before they were run (Appendix~\ref{app:seed}). Everything reported per block is produced by one analysis pass from the packaged outcomes; the leaderboard, the write-flow share, and the per-image examples are descriptive summaries of those outputs added after the grid was complete.

\paragraph{Artifact.} The release contains the evaluation code and adapters, the registered protocol file, the cohort manifests with their seeds, the fitted directions, the per-image outcomes of every module (write, dose, refit, locus, and template rows), every block's scorecard and coverage file, and the rendered prompts; a new checkpoint is added with one adapter that names its consumed block and passes the gates.
